\scanpage{40}
Soient \inserted{$k$ anneau comm.} $\unsure{\mathcal C|\mathcal E|\mathbb K}$ une catégorie \inserted{$k$-}additive (i.e. additive, munie
de $k\to\operatorname{End}(\mathrm{id}_{\mathcal C})$), $C^\bullet$ un complexe \inserted{de cochaînes} dans
$\mathcal C$
\[
 \cdots\longrightarrow C^0\xrightarrow{d^0}C^1\xrightarrow{d^1}
 \cdots\longrightarrow C^i\xrightarrow{d^i}C^{i+1}\longrightarrow\cdots
\]
On suppose :

a)
\[
\operatorname{Hom}(C^i,C^j)\simeq
\begin{cases}
0&\text{si }j\ne i,i+1\\
k&\text{si }j=i,\ \text{base }\mathrm{id}_{C^i}\\
k&\text{si }j=i+1,\ \text{base }d^i
\end{cases}
\]

b) Tout objet de $\mathcal C$ est isomorphe à une
somme \inserted{finie} d'objets de la forme $C^i$.

Soit $\mathcal A$ la sous-catégorie pleine de $\mathcal C$
formée des \deleted{$\unsure{\Phi}$} $C^i$ ; elle n'est pas additive,
mais hérite d'une structure $k$-linéaire,
i.e les $\operatorname{Hom}_{\mathcal A}$ sont des $k$-modules, leur
composition est $k$-bil.

\deleted{\unsure{C|E}} L'existence de $(\mathcal C,C^\bullet)$ satisfaisant a) et
b) est immédiate par construction ``abstraite''.
On peut aussi réaliser $\mathcal C$ comme \inserted{la} cat. \inserted{$K_{\bullet}$} des
complexes de chaînes \inserted{$L_{\bullet}$} de $k$-modules
satisfaisant les conditions

$1^\circ$) \deleted{\unsure{$K_{\bullet}$}} $L_{\bullet}$ est acyclique i.e. $H_i(L_{\bullet})=0$ $\forall i\geq0$

$2^\circ$) $L_n$ et $Z_n(L_{\bullet})$ sont libres de t. f. pour tt $n$

$3^\circ$) $L_n$ est nul pour $n$ grand.

On prendra alors :
\[
 C^i=\deleted{\unsure{K^i_{\bullet}}} C^i_{\bullet}=
 \{\,0\longrightarrow\cdots\longrightarrow0\longrightarrow
 \overset{\text{deg }i+1}{k}\xrightarrow{\mathrm{id}}
 \overset{\text{deg }i}{k}\longrightarrow0\longrightarrow\cdots\longrightarrow
 \overset{\text{deg }0}{0}\,\}
\]
i.e. les comp. sont nulles sauf en degré $i,i+1$ où on
aura $k$, et l'app. diff. \deleted{\unsure{non nulle}}\inserted{$i+1$} est l'identité.
On prend $d^i\colon \deleted{\unsure{K}}C^i_{\bullet}\xrightarrow{d^i_{\bullet}}
\deleted{\unsure{K}}C^{i+1}_{\bullet}$ caractérisé
par $(d^i)_{i+1}=\mathrm{id}_k$. Les conditions a) et b) sont
satisfaites trivialement. On pourrait aussi
prendre des complexes de cochaînes à la place,
% margin note written sideways
\margin{avec les $C_i^\bullet=(\overset{\text{deg }0}{0}\to\cdots\to0\to
\overset{\text{deg }i}{k}\to\overset{\text{deg }i+1}{k}\to\cdots$, mais
ici on a $C_i^\bullet\longleftarrow C_{i+1}^\bullet$, \deleted{\unsure{ce}} il faut en prendre
la cat. \uline{opposée} de celle \inserted{des} cochaînes \unsure{en quest.|envisagée}.
Mais on le retrouvera en \unsure{2 repr.|? repr.|2 vers.} de façon plus
\uline{élégante} dans les Ex.}
% archive sheet number: 39
